\section{Введение}

Кластеризация является важным инструментом анализа данных, который находит широкое применение во многих областях, включая биоинформатику, медицину, финансы, исследования социальных сетей и многие другие. Однако с увеличением размерности данных стандартные методы кластеризации могут столкнуться с проблемой низкого качества разбиения.

Проблема качества кластеризации особенно остра в многомерном пространстве, где количество признаков значительно превышает количество объектов. Это явление, известное как <<проклятие размерности>>, приводит к тому, что традиционные методы кластеризации становятся менее эффективными из-за разреженности данных и трудностей в интерпретации результатов.

Одним из способов преодоления этой проблемы является привлечение учителя для предоставления дополнительной информации о данных. Однако разметка больших наборов данных для обучения моделей машинного обучения с учителем требует значительных временных и финансовых затрат. В связи с этим возникает потребность в разработке методов кластеризации, которые могли бы использовать только часть размеченных данных для улучшения качества разбиения.

В данной работе исследуется подход кластеризации с частичным привлечением учителя (\ssClustering{}), который позволяет использовать ограниченное количество размеченных данных для повышения эффективности кластеризации в пространствах большой размерности.

Для достижения данной цели поставлены следующие задачи:
\begin{enumerate}
\item систематизировать и проанализировать существующие методы кластеризации,
\item реализовать алгоритм кластеризации с частичным привлечением учителя в пространствах большой размерности,
\item провести экспериментальное исследование и анализ влияния параметров алгоритма на качество кластеризации.
\end{enumerate}

В результате выполнения поставленных задач будет разработан новый метод кластеризации с частичным привлечением учителя. Особенностью предложенного метода является его способность эффективно работать с ограниченным объемом размеченных данных, что делает его особенно привлекательным для сфер, где разметка данных является сложной, трудоемкой и дорогостоящей задачей, включая анализ текста, аудиосигналов изображений и других.
Кроме того, данная работа предоставит возможность глубже исследовать и понять, как ограниченный объем размеченных данных может быть использован для улучшения качества кластеризации. Это позволит существенно оптимизировать процесс разметки данных и снизить затраты на него.
